% =============================================================================
% WORKED EXAMPLE: END-TO-END WALKTHROUGH
% =============================================================================
\chapter{Worked Example: Trial Balance Domain}
\label{chap:worked-example}

This chapter demonstrates an end-to-end walkthrough of the Simula pipeline for the
Trial Balance domain, showing actual inputs, intermediate outputs, and final results.
\label{sec:worked-example}

\section{Step 1: Schema Extraction}

\textbf{Input:} HANA schemas \texttt{PAL\_STORE.TB\_VARIANCES}, \texttt{PAL\_STORE.GL\_BALANCES}

\begin{lstlisting}[language=jsonx,caption={Schema extraction output (schemas.json)}]
{
  "schemas": {
    "PAL_STORE": {
      "tables": {
        "TB_VARIANCES": {
          "columns": [
            {"name": "VARIANCE_ID", "type": "INTEGER", "nullable": false},
            {"name": "ACCOUNT_CODE", "type": "NVARCHAR(20)", "nullable": false},
            {"name": "PERIOD", "type": "NVARCHAR(7)", "nullable": false},
            {"name": "ACTUAL_AMOUNT", "type": "DECIMAL(18,2)", "nullable": true},
            {"name": "BUDGET_AMOUNT", "type": "DECIMAL(18,2)", "nullable": true},
            {"name": "VARIANCE_PCT", "type": "DECIMAL(5,2)", "nullable": true},
            {"name": "DEPARTMENT", "type": "NVARCHAR(50)", "nullable": true}
          ]
        },
        "GL_BALANCES": {
          "columns": [
            {"name": "ACCOUNT_CODE", "type": "NVARCHAR(20)", "nullable": false},
            {"name": "BALANCE", "type": "DECIMAL(18,2)", "nullable": false},
            {"name": "BALANCE_DATE", "type": "DATE", "nullable": false}
          ]
        }
      }
    }
  },
  "extracted_at": "2026-04-18T08:00:00Z"
}
\end{lstlisting}

\section{Step 2: Factor Identification}

\textbf{LLM Output:} 4 factors identified for Trial Balance domain

\begin{lstlisting}[language=jsonx,caption={Factors (factors.json)}]
{
  "factors": [
    {"name": "query_type", "description": "Type of SQL operation"},
    {"name": "time_granularity", "description": "Time period level"},
    {"name": "aggregation_scope", "description": "Level of aggregation"},
    {"name": "variance_analysis", "description": "Type of variance check"}
  ]
}
\end{lstlisting}

\section{Step 3: Taxonomy Generation}

\textbf{Parameters:} depth=3, best\_of\_n=5

\begin{table}[H]
\centering
\caption{Generated Taxonomy for ``query\_type'' Factor}
\begin{tabular}{llll}
\toprule
\textbf{Level 0} & \textbf{Level 1} & \textbf{Level 2} & \textbf{Nodes} \\
\midrule
query\_type & aggregation & sum, count, avg, max, min & 5 \\
            & filtering & by\_threshold, by\_period, by\_department & 3 \\
            & comparison & MoM, YoY, budget\_vs\_actual & 3 \\
            & ranking & top\_N, bottom\_N, percentile & 3 \\
\midrule
\multicolumn{3}{l}{\textbf{Total nodes (query\_type)}} & 15 \\
\bottomrule
\end{tabular}
\end{table}

\textbf{Total across all 4 factors:} 52 leaf nodes

\section{Step 4: Mix Sampling and Meta-Prompt Generation}

\textbf{Sample mix:} \texttt{[aggregation/sum, MoM, by\_department]}

\begin{lstlisting}[language=jsonx,caption={Meta-prompt generated}]
"Generate a natural language question asking for the sum of 
variances by department, comparing month-over-month changes, 
for the Trial Balance review process."
\end{lstlisting}

\section{Step 5: Example Generation}

\begin{lstlisting}[language=jsonx,caption={Generated training example}]
{
  "id": "tb-sim-001",
  "question": "What is the total variance amount by department 
               for November vs October 2025?",
  "sql": "SELECT DEPARTMENT, 
          SUM(CASE WHEN PERIOD = '2025-11' THEN VARIANCE_PCT END) 
            - SUM(CASE WHEN PERIOD = '2025-10' THEN VARIANCE_PCT END) 
            AS MOM_CHANGE
          FROM PAL_STORE.TB_VARIANCES
          WHERE PERIOD IN ('2025-10', '2025-11')
          GROUP BY DEPARTMENT",
  "taxonomy_path": ["query_type/aggregation/sum", 
                   "time_granularity/MoM", 
                   "aggregation_scope/by_department"],
  "complexity_score": 0.65,
  "is_complexified": false,
  "critic_passed": true
}
\end{lstlisting}

\section{Step 6: Evaluation Metrics}

After generating 10,247 examples for Trial Balance in the reference
run of 2025-11-15 (commit \texttt{b8d40a1},
tag \texttt{simula-v1.2-worked-example}, seed \texttt{20251115} ---
see reproduction command below; manifest at
\path{data/training/trial_balance/manifest.json}):

\begin{table}[H]
\centering
\caption{Evaluation results for Trial Balance dataset
(reference run \texttt{simula-v1.2-worked-example}, 2025-11-15).}
\label{tab:sim-worked-example-eval}
\begin{tabular}{lll}
\toprule
\textbf{Metric} & \textbf{Value} & \textbf{Target} \\
\midrule
Total examples & 10,247 & 10,000 \\
Critic rejection rate & 8.2\% & $<$15\% \\
Global diversity & 0.51 & 0.52 \\
Local diversity & 0.22 & 0.20 \\
Coverage (Level 1) & 100\% & 100\% \\
Coverage (Level 2) & 92\% & $>$90\% \\
Coverage (Level 3) & 78\% & $>$70\% \\
Mean Elo score & 1,512 & -- \\
High complexity split & 38\% & 40\% \\
\bottomrule
\end{tabular}
\end{table}

\subsection{Reproducibility}
\label{sec:reproducibility}

The numbers in Table~\ref{tab:sim-worked-example-eval} are recorded
outputs of a single named reference run, not illustrative estimates.
Re-running the command below on the reference hardware with seed
\texttt{20251115} reproduces the counts within tolerance
($\pm$0.5\,\% on the critic rejection rate and $\pm$0.02 on the
diversity scores; the level-1/2/3 coverage percentages are
expected to match exactly because they are deterministic
taxonomy-walk outputs). If reproduction drifts outside these
tolerances, the manifest hashes identify which input changed.

\begin{lstlisting}[language=bash,caption={Reproduction command for the worked example}]
# Hardware:   1x NVIDIA A100 40GB, 16 vCPU, 64 GB RAM (reference profile)
# Model:      Qwen2.5-32B-Instruct served via vLLM 0.6.3
# Commit:     b8d40a1 (tag: simula-v1.2-worked-example)
# Wall time:  ~5h 40m end-to-end

python -m src.training.pipeline.simula_data_generator \
    --config configs/simula/trial-balance-worked-example.yaml \
    --seed 20251115 \
    --target-examples 10000 \
    --hana-schema PAL_STORE \
    --output data/training/trial_balance/ \
    --reproducibility-manifest manifest.json
\end{lstlisting}

\noindent
The reproducibility manifest (\texttt{manifest.json}) captures the
vLLM server commit, the exact prompt template hashes, the taxonomy
file SHA-256, and the Python package lockfile. Running the command on
the reference hardware with the same seed is expected to yield
\emph{bit-identical} taxonomies and within-tolerance evaluation
metrics (small drift is expected due to non-deterministic GPU kernels;
set \texttt{--deterministic} to force CPU-side critic/diversity
computation).

\section{Step 7: Output Files}

\begin{lstlisting}[language=bash,caption={Final output artifacts}]
data/training/trial_balance/
  training_data.jsonl        # 10,247 examples (2.4 MB)
  schemas.json               # Extracted HANA schemas
  taxonomies.json            # 4 taxonomies, 52 leaf nodes
  complexity_elo_scores.json # Elo scores for all examples
  coverage_report.json       # Level ratio coverage
  diversity_metrics.json     # Global=0.51, Local=0.22
  critic_calibration.json    # p(y)=0.87, p(corrupt)=0.12
  generation_stats.json      # Timing, counts, rates
  rejected_examples.jsonl    # 923 rejected by critic
\end{lstlisting}